\documentclass{article}

% ============================================================
% INFR11269 Machine Learning Systems -- Coursework Report (skeleton)
% Built from report/latex_template/report.tex. The .sty/.bst files in this
% directory are unmodified copies of report/latex_template/.
%
% Format rules kept from the template: no geometry/font/spacing changes,
% \iclrfinalcopy stays commented out, abstract environment stays commented
% out, \author and \lhead lines unchanged, References before the Appendix.
%
% Numbers drop in from rag/extensions/analysis/summarise.py:
%   python rag/extensions/analysis/summarise.py --split dev      --out report/asr-redundancy
%   python rag/extensions/analysis/summarise.py --split practice --out report/asr-redundancy
% which writes tables/asr_{split}.tex and figures/{breakdown,tradeoff,speedup}_{split}.pdf
% next to this file. Until then every slot renders as a red \TBD box.
% ============================================================

\usepackage{iclr2027_conference,times}

\usepackage{hyperref}
\usepackage{url}
\usepackage{booktabs}
\usepackage{graphicx}
\usepackage{amsmath}
\usepackage{amssymb}
\usepackage{xcolor}
\usepackage{verbatim}

% The official ICLR archive includes math_commands.tex. It is optional:
%%%%% NEW MATH DEFINITIONS %%%%%

\usepackage{amsmath,amsfonts,bm}

% Mark sections of captions for referring to divisions of figures
\newcommand{\figleft}{{\em (Left)}}
\newcommand{\figcenter}{{\em (Center)}}
\newcommand{\figright}{{\em (Right)}}
\newcommand{\figtop}{{\em (Top)}}
\newcommand{\figbottom}{{\em (Bottom)}}
\newcommand{\captiona}{{\em (a)}}
\newcommand{\captionb}{{\em (b)}}
\newcommand{\captionc}{{\em (c)}}
\newcommand{\captiond}{{\em (d)}}

% Highlight a newly defined term
\newcommand{\newterm}[1]{{\bf #1}}


% Figure reference, lower-case.
\def\figref#1{figure~\ref{#1}}
% Figure reference, capital. For start of sentence
\def\Figref#1{Figure~\ref{#1}}
\def\twofigref#1#2{figures \ref{#1} and \ref{#2}}
\def\quadfigref#1#2#3#4{figures \ref{#1}, \ref{#2}, \ref{#3} and \ref{#4}}
% Section reference, lower-case.
\def\secref#1{section~\ref{#1}}
% Section reference, capital.
\def\Secref#1{Section~\ref{#1}}
% Reference to two sections.
\def\twosecrefs#1#2{sections \ref{#1} and \ref{#2}}
% Reference to three sections.
\def\secrefs#1#2#3{sections \ref{#1}, \ref{#2} and \ref{#3}}
% Reference to an equation, lower-case.
\def\eqref#1{equation~\ref{#1}}
% Reference to an equation, upper case
\def\Eqref#1{Equation~\ref{#1}}
% A raw reference to an equation---avoid using if possible
\def\plaineqref#1{\ref{#1}}
% Reference to a chapter, lower-case.
\def\chapref#1{chapter~\ref{#1}}
% Reference to an equation, upper case.
\def\Chapref#1{Chapter~\ref{#1}}
% Reference to a range of chapters
\def\rangechapref#1#2{chapters\ref{#1}--\ref{#2}}
% Reference to an algorithm, lower-case.
\def\algref#1{algorithm~\ref{#1}}
% Reference to an algorithm, upper case.
\def\Algref#1{Algorithm~\ref{#1}}
\def\twoalgref#1#2{algorithms \ref{#1} and \ref{#2}}
\def\Twoalgref#1#2{Algorithms \ref{#1} and \ref{#2}}
% Reference to a part, lower case
\def\partref#1{part~\ref{#1}}
% Reference to a part, upper case
\def\Partref#1{Part~\ref{#1}}
\def\twopartref#1#2{parts \ref{#1} and \ref{#2}}

\def\ceil#1{\lceil #1 \rceil}
\def\floor#1{\lfloor #1 \rfloor}
\def\1{\bm{1}}
\newcommand{\train}{\mathcal{D}}
\newcommand{\valid}{\mathcal{D_{\mathrm{valid}}}}
\newcommand{\test}{\mathcal{D_{\mathrm{test}}}}

\def\eps{{\epsilon}}


% Random variables
\def\reta{{\textnormal{$\eta$}}}
\def\ra{{\textnormal{a}}}
\def\rb{{\textnormal{b}}}
\def\rc{{\textnormal{c}}}
\def\rd{{\textnormal{d}}}
\def\re{{\textnormal{e}}}
\def\rf{{\textnormal{f}}}
\def\rg{{\textnormal{g}}}
\def\rh{{\textnormal{h}}}
\def\ri{{\textnormal{i}}}
\def\rj{{\textnormal{j}}}
\def\rk{{\textnormal{k}}}
\def\rl{{\textnormal{l}}}
% rm is already a command, just don't name any random variables m
\def\rn{{\textnormal{n}}}
\def\ro{{\textnormal{o}}}
\def\rp{{\textnormal{p}}}
\def\rq{{\textnormal{q}}}
\def\rr{{\textnormal{r}}}
\def\rs{{\textnormal{s}}}
\def\rt{{\textnormal{t}}}
\def\ru{{\textnormal{u}}}
\def\rv{{\textnormal{v}}}
\def\rw{{\textnormal{w}}}
\def\rx{{\textnormal{x}}}
\def\ry{{\textnormal{y}}}
\def\rz{{\textnormal{z}}}

% Random vectors
\def\rvepsilon{{\mathbf{\epsilon}}}
\def\rvtheta{{\mathbf{\theta}}}
\def\rva{{\mathbf{a}}}
\def\rvb{{\mathbf{b}}}
\def\rvc{{\mathbf{c}}}
\def\rvd{{\mathbf{d}}}
\def\rve{{\mathbf{e}}}
\def\rvf{{\mathbf{f}}}
\def\rvg{{\mathbf{g}}}
\def\rvh{{\mathbf{h}}}
\def\rvu{{\mathbf{i}}}
\def\rvj{{\mathbf{j}}}
\def\rvk{{\mathbf{k}}}
\def\rvl{{\mathbf{l}}}
\def\rvm{{\mathbf{m}}}
\def\rvn{{\mathbf{n}}}
\def\rvo{{\mathbf{o}}}
\def\rvp{{\mathbf{p}}}
\def\rvq{{\mathbf{q}}}
\def\rvr{{\mathbf{r}}}
\def\rvs{{\mathbf{s}}}
\def\rvt{{\mathbf{t}}}
\def\rvu{{\mathbf{u}}}
\def\rvv{{\mathbf{v}}}
\def\rvw{{\mathbf{w}}}
\def\rvx{{\mathbf{x}}}
\def\rvy{{\mathbf{y}}}
\def\rvz{{\mathbf{z}}}

% Elements of random vectors
\def\erva{{\textnormal{a}}}
\def\ervb{{\textnormal{b}}}
\def\ervc{{\textnormal{c}}}
\def\ervd{{\textnormal{d}}}
\def\erve{{\textnormal{e}}}
\def\ervf{{\textnormal{f}}}
\def\ervg{{\textnormal{g}}}
\def\ervh{{\textnormal{h}}}
\def\ervi{{\textnormal{i}}}
\def\ervj{{\textnormal{j}}}
\def\ervk{{\textnormal{k}}}
\def\ervl{{\textnormal{l}}}
\def\ervm{{\textnormal{m}}}
\def\ervn{{\textnormal{n}}}
\def\ervo{{\textnormal{o}}}
\def\ervp{{\textnormal{p}}}
\def\ervq{{\textnormal{q}}}
\def\ervr{{\textnormal{r}}}
\def\ervs{{\textnormal{s}}}
\def\ervt{{\textnormal{t}}}
\def\ervu{{\textnormal{u}}}
\def\ervv{{\textnormal{v}}}
\def\ervw{{\textnormal{w}}}
\def\ervx{{\textnormal{x}}}
\def\ervy{{\textnormal{y}}}
\def\ervz{{\textnormal{z}}}

% Random matrices
\def\rmA{{\mathbf{A}}}
\def\rmB{{\mathbf{B}}}
\def\rmC{{\mathbf{C}}}
\def\rmD{{\mathbf{D}}}
\def\rmE{{\mathbf{E}}}
\def\rmF{{\mathbf{F}}}
\def\rmG{{\mathbf{G}}}
\def\rmH{{\mathbf{H}}}
\def\rmI{{\mathbf{I}}}
\def\rmJ{{\mathbf{J}}}
\def\rmK{{\mathbf{K}}}
\def\rmL{{\mathbf{L}}}
\def\rmM{{\mathbf{M}}}
\def\rmN{{\mathbf{N}}}
\def\rmO{{\mathbf{O}}}
\def\rmP{{\mathbf{P}}}
\def\rmQ{{\mathbf{Q}}}
\def\rmR{{\mathbf{R}}}
\def\rmS{{\mathbf{S}}}
\def\rmT{{\mathbf{T}}}
\def\rmU{{\mathbf{U}}}
\def\rmV{{\mathbf{V}}}
\def\rmW{{\mathbf{W}}}
\def\rmX{{\mathbf{X}}}
\def\rmY{{\mathbf{Y}}}
\def\rmZ{{\mathbf{Z}}}

% Elements of random matrices
\def\ermA{{\textnormal{A}}}
\def\ermB{{\textnormal{B}}}
\def\ermC{{\textnormal{C}}}
\def\ermD{{\textnormal{D}}}
\def\ermE{{\textnormal{E}}}
\def\ermF{{\textnormal{F}}}
\def\ermG{{\textnormal{G}}}
\def\ermH{{\textnormal{H}}}
\def\ermI{{\textnormal{I}}}
\def\ermJ{{\textnormal{J}}}
\def\ermK{{\textnormal{K}}}
\def\ermL{{\textnormal{L}}}
\def\ermM{{\textnormal{M}}}
\def\ermN{{\textnormal{N}}}
\def\ermO{{\textnormal{O}}}
\def\ermP{{\textnormal{P}}}
\def\ermQ{{\textnormal{Q}}}
\def\ermR{{\textnormal{R}}}
\def\ermS{{\textnormal{S}}}
\def\ermT{{\textnormal{T}}}
\def\ermU{{\textnormal{U}}}
\def\ermV{{\textnormal{V}}}
\def\ermW{{\textnormal{W}}}
\def\ermX{{\textnormal{X}}}
\def\ermY{{\textnormal{Y}}}
\def\ermZ{{\textnormal{Z}}}

% Vectors
\def\vzero{{\bm{0}}}
\def\vone{{\bm{1}}}
\def\vmu{{\bm{\mu}}}
\def\vtheta{{\bm{\theta}}}
\def\va{{\bm{a}}}
\def\vb{{\bm{b}}}
\def\vc{{\bm{c}}}
\def\vd{{\bm{d}}}
\def\ve{{\bm{e}}}
\def\vf{{\bm{f}}}
\def\vg{{\bm{g}}}
\def\vh{{\bm{h}}}
\def\vi{{\bm{i}}}
\def\vj{{\bm{j}}}
\def\vk{{\bm{k}}}
\def\vl{{\bm{l}}}
\def\vm{{\bm{m}}}
\def\vn{{\bm{n}}}
\def\vo{{\bm{o}}}
\def\vp{{\bm{p}}}
\def\vq{{\bm{q}}}
\def\vr{{\bm{r}}}
\def\vs{{\bm{s}}}
\def\vt{{\bm{t}}}
\def\vu{{\bm{u}}}
\def\vv{{\bm{v}}}
\def\vw{{\bm{w}}}
\def\vx{{\bm{x}}}
\def\vy{{\bm{y}}}
\def\vz{{\bm{z}}}

% Elements of vectors
\def\evalpha{{\alpha}}
\def\evbeta{{\beta}}
\def\evepsilon{{\epsilon}}
\def\evlambda{{\lambda}}
\def\evomega{{\omega}}
\def\evmu{{\mu}}
\def\evpsi{{\psi}}
\def\evsigma{{\sigma}}
\def\evtheta{{\theta}}
\def\eva{{a}}
\def\evb{{b}}
\def\evc{{c}}
\def\evd{{d}}
\def\eve{{e}}
\def\evf{{f}}
\def\evg{{g}}
\def\evh{{h}}
\def\evi{{i}}
\def\evj{{j}}
\def\evk{{k}}
\def\evl{{l}}
\def\evm{{m}}
\def\evn{{n}}
\def\evo{{o}}
\def\evp{{p}}
\def\evq{{q}}
\def\evr{{r}}
\def\evs{{s}}
\def\evt{{t}}
\def\evu{{u}}
\def\evv{{v}}
\def\evw{{w}}
\def\evx{{x}}
\def\evy{{y}}
\def\evz{{z}}

% Matrix
\def\mA{{\bm{A}}}
\def\mB{{\bm{B}}}
\def\mC{{\bm{C}}}
\def\mD{{\bm{D}}}
\def\mE{{\bm{E}}}
\def\mF{{\bm{F}}}
\def\mG{{\bm{G}}}
\def\mH{{\bm{H}}}
\def\mI{{\bm{I}}}
\def\mJ{{\bm{J}}}
\def\mK{{\bm{K}}}
\def\mL{{\bm{L}}}
\def\mM{{\bm{M}}}
\def\mN{{\bm{N}}}
\def\mO{{\bm{O}}}
\def\mP{{\bm{P}}}
\def\mQ{{\bm{Q}}}
\def\mR{{\bm{R}}}
\def\mS{{\bm{S}}}
\def\mT{{\bm{T}}}
\def\mU{{\bm{U}}}
\def\mV{{\bm{V}}}
\def\mW{{\bm{W}}}
\def\mX{{\bm{X}}}
\def\mY{{\bm{Y}}}
\def\mZ{{\bm{Z}}}
\def\mBeta{{\bm{\beta}}}
\def\mPhi{{\bm{\Phi}}}
\def\mLambda{{\bm{\Lambda}}}
\def\mSigma{{\bm{\Sigma}}}

% Tensor
\DeclareMathAlphabet{\mathsfit}{\encodingdefault}{\sfdefault}{m}{sl}
\SetMathAlphabet{\mathsfit}{bold}{\encodingdefault}{\sfdefault}{bx}{n}
\newcommand{\tens}[1]{\bm{\mathsfit{#1}}}
\def\tA{{\tens{A}}}
\def\tB{{\tens{B}}}
\def\tC{{\tens{C}}}
\def\tD{{\tens{D}}}
\def\tE{{\tens{E}}}
\def\tF{{\tens{F}}}
\def\tG{{\tens{G}}}
\def\tH{{\tens{H}}}
\def\tI{{\tens{I}}}
\def\tJ{{\tens{J}}}
\def\tK{{\tens{K}}}
\def\tL{{\tens{L}}}
\def\tM{{\tens{M}}}
\def\tN{{\tens{N}}}
\def\tO{{\tens{O}}}
\def\tP{{\tens{P}}}
\def\tQ{{\tens{Q}}}
\def\tR{{\tens{R}}}
\def\tS{{\tens{S}}}
\def\tT{{\tens{T}}}
\def\tU{{\tens{U}}}
\def\tV{{\tens{V}}}
\def\tW{{\tens{W}}}
\def\tX{{\tens{X}}}
\def\tY{{\tens{Y}}}
\def\tZ{{\tens{Z}}}


% Graph
\def\gA{{\mathcal{A}}}
\def\gB{{\mathcal{B}}}
\def\gC{{\mathcal{C}}}
\def\gD{{\mathcal{D}}}
\def\gE{{\mathcal{E}}}
\def\gF{{\mathcal{F}}}
\def\gG{{\mathcal{G}}}
\def\gH{{\mathcal{H}}}
\def\gI{{\mathcal{I}}}
\def\gJ{{\mathcal{J}}}
\def\gK{{\mathcal{K}}}
\def\gL{{\mathcal{L}}}
\def\gM{{\mathcal{M}}}
\def\gN{{\mathcal{N}}}
\def\gO{{\mathcal{O}}}
\def\gP{{\mathcal{P}}}
\def\gQ{{\mathcal{Q}}}
\def\gR{{\mathcal{R}}}
\def\gS{{\mathcal{S}}}
\def\gT{{\mathcal{T}}}
\def\gU{{\mathcal{U}}}
\def\gV{{\mathcal{V}}}
\def\gW{{\mathcal{W}}}
\def\gX{{\mathcal{X}}}
\def\gY{{\mathcal{Y}}}
\def\gZ{{\mathcal{Z}}}

% Sets
\def\sA{{\mathbb{A}}}
\def\sB{{\mathbb{B}}}
\def\sC{{\mathbb{C}}}
\def\sD{{\mathbb{D}}}
% Don't use a set called E, because this would be the same as our symbol
% for expectation.
\def\sF{{\mathbb{F}}}
\def\sG{{\mathbb{G}}}
\def\sH{{\mathbb{H}}}
\def\sI{{\mathbb{I}}}
\def\sJ{{\mathbb{J}}}
\def\sK{{\mathbb{K}}}
\def\sL{{\mathbb{L}}}
\def\sM{{\mathbb{M}}}
\def\sN{{\mathbb{N}}}
\def\sO{{\mathbb{O}}}
\def\sP{{\mathbb{P}}}
\def\sQ{{\mathbb{Q}}}
\def\sR{{\mathbb{R}}}
\def\sS{{\mathbb{S}}}
\def\sT{{\mathbb{T}}}
\def\sU{{\mathbb{U}}}
\def\sV{{\mathbb{V}}}
\def\sW{{\mathbb{W}}}
\def\sX{{\mathbb{X}}}
\def\sY{{\mathbb{Y}}}
\def\sZ{{\mathbb{Z}}}

% Entries of a matrix
\def\emLambda{{\Lambda}}
\def\emA{{A}}
\def\emB{{B}}
\def\emC{{C}}
\def\emD{{D}}
\def\emE{{E}}
\def\emF{{F}}
\def\emG{{G}}
\def\emH{{H}}
\def\emI{{I}}
\def\emJ{{J}}
\def\emK{{K}}
\def\emL{{L}}
\def\emM{{M}}
\def\emN{{N}}
\def\emO{{O}}
\def\emP{{P}}
\def\emQ{{Q}}
\def\emR{{R}}
\def\emS{{S}}
\def\emT{{T}}
\def\emU{{U}}
\def\emV{{V}}
\def\emW{{W}}
\def\emX{{X}}
\def\emY{{Y}}
\def\emZ{{Z}}
\def\emSigma{{\Sigma}}

% entries of a tensor
% Same font as tensor, without \bm wrapper
\newcommand{\etens}[1]{\mathsfit{#1}}
\def\etLambda{{\etens{\Lambda}}}
\def\etA{{\etens{A}}}
\def\etB{{\etens{B}}}
\def\etC{{\etens{C}}}
\def\etD{{\etens{D}}}
\def\etE{{\etens{E}}}
\def\etF{{\etens{F}}}
\def\etG{{\etens{G}}}
\def\etH{{\etens{H}}}
\def\etI{{\etens{I}}}
\def\etJ{{\etens{J}}}
\def\etK{{\etens{K}}}
\def\etL{{\etens{L}}}
\def\etM{{\etens{M}}}
\def\etN{{\etens{N}}}
\def\etO{{\etens{O}}}
\def\etP{{\etens{P}}}
\def\etQ{{\etens{Q}}}
\def\etR{{\etens{R}}}
\def\etS{{\etens{S}}}
\def\etT{{\etens{T}}}
\def\etU{{\etens{U}}}
\def\etV{{\etens{V}}}
\def\etW{{\etens{W}}}
\def\etX{{\etens{X}}}
\def\etY{{\etens{Y}}}
\def\etZ{{\etens{Z}}}

% The true underlying data generating distribution
\newcommand{\pdata}{p_{\rm{data}}}
% The empirical distribution defined by the training set
\newcommand{\ptrain}{\hat{p}_{\rm{data}}}
\newcommand{\Ptrain}{\hat{P}_{\rm{data}}}
% The model distribution
\newcommand{\pmodel}{p_{\rm{model}}}
\newcommand{\Pmodel}{P_{\rm{model}}}
\newcommand{\ptildemodel}{\tilde{p}_{\rm{model}}}
% Stochastic autoencoder distributions
\newcommand{\pencode}{p_{\rm{encoder}}}
\newcommand{\pdecode}{p_{\rm{decoder}}}
\newcommand{\precons}{p_{\rm{reconstruct}}}

\newcommand{\laplace}{\mathrm{Laplace}} % Laplace distribution

\newcommand{\E}{\mathbb{E}}
\newcommand{\Ls}{\mathcal{L}}
\newcommand{\R}{\mathbb{R}}
\newcommand{\emp}{\tilde{p}}
\newcommand{\lr}{\alpha}
\newcommand{\reg}{\lambda}
\newcommand{\rect}{\mathrm{rectifier}}
\newcommand{\softmax}{\mathrm{softmax}}
\newcommand{\sigmoid}{\sigma}
\newcommand{\softplus}{\zeta}
\newcommand{\KL}{D_{\mathrm{KL}}}
\newcommand{\Var}{\mathrm{Var}}
\newcommand{\standarderror}{\mathrm{SE}}
\newcommand{\Cov}{\mathrm{Cov}}
% Wolfram Mathworld says $L^2$ is for function spaces and $\ell^2$ is for vectors
% But then they seem to use $L^2$ for vectors throughout the site, and so does
% wikipedia.
\newcommand{\normlzero}{L^0}
\newcommand{\normlone}{L^1}
\newcommand{\normltwo}{L^2}
\newcommand{\normlp}{L^p}
\newcommand{\normmax}{L^\infty}

\newcommand{\parents}{Pa} % See usage in notation.tex. Chosen to match Daphne's book.

\DeclareMathOperator*{\argmax}{arg\,max}
\DeclareMathOperator*{\argmin}{arg\,min}

\DeclareMathOperator{\sign}{sign}
\DeclareMathOperator{\Tr}{Tr}
\let\ab\allowbreak


% ---------- placeholder macros (remove nothing silently) ----------
\newcommand{\TBD}[1]{\textcolor{red}{\textbf{[TBD: #1]}}}
\newcommand{\NOTE}[1]{\textcolor{blue}{\textit{[Note: #1]}}}
% \inputorTBD{file}{what goes here}
\newcommand{\inputorTBD}[2]{\IfFileExists{#1}{\input{#1}}{\TBD{#2 (\texttt{\detokenize{#1}} not generated yet)}}}
% \figorTBD{file}{width}{what goes here}
\newcommand{\figorTBD}[3]{\IfFileExists{#1}{\includegraphics[width=#2]{#1}}{\fbox{\parbox[c][3.2cm][c]{#2}{\centering\TBD{#3}\\[2pt]\scriptsize\texttt{\detokenize{#1}}}}}}
% \logorTBD{file}{what goes here}
\newcommand{\logorTBD}[2]{\IfFileExists{#1}{{\footnotesize\verbatiminput{#1}}}{\TBD{#2 (\texttt{\detokenize{#1}} missing)}}}

\title{Removing Redundant Computation from a\\Speech-to-Text Front End}

% Keep this line unchanged.
\author{Anonymous coursework submission}

% Keep this COMMENTED OUT.
% Review mode is part of the coursework format and is checked automatically.
% \iclrfinalcopy

\begin{document}

\maketitle

% Course-specific template fingerprint. Keep this line unchanged.
\lhead{Under review as a report for UoE 26/27 Sem1 INFR11269: Machine Learning Systems}

% Keep this COMMENTED OUT.  (The template forbids the abstract environment;
% the draft abstract is kept here as a comment so it can be reused in the
% introduction or the video.)
% \begin{abstract}
% \end{abstract}
%
% DRAFT ABSTRACT (4 sentences, not rendered):
% In an audio-query RAG pipeline, the speech-to-text stage spends most of its
% time on work whose result is already known: re-running the decoder over the
% whole prefix for every token, and encoding up to 30 s of padding for queries
% that average about 5 s.
% We remove the first redundancy exactly with a KV cache built on our own CUDA
% kernels, and the second approximately by trimming the encoder input to the
% real audio length plus a margin.
% On a pinned RTX 2080 Ti, the KV cache cuts mean ASR latency 2.0x with
% bit-identical transcripts, and trimming to the real length brings the total to
% 4.8x with no Recall@10 change distinguishable from zero (dev, n=125).
% A cost model written before measurement predicts the encoder savings within
% 7%, and explains why decoding then becomes a weight-bandwidth-bound GEMV.

% ============================================================
\section{Introduction}
\label{sec:introduction}

The system under study answers spoken questions: a 16\,kHz audio query is
transcribed by an automatic speech recognition (ASR) model, the transcript is
embedded and used for dense retrieval over a 24{,}000-passage corpus, and the
retrieved passages are passed to a hosted language model (the ELM) that writes the
answer \citep{lewis2020rag}. Retrieval only ever sees the transcript, so the ASR
stage is both the first latency cost and the only source of input error for
everything downstream.

The supplied ASR model, GLM-ASR-Nano, runs on CUDA kernels we wrote for HW1. Its
generation loop has two sources of redundant work. First, it has no key--value
(KV) cache: each new token re-runs all 28 decoder layers over the entire prefix.
Second, every clip is padded to 30\,s as in Whisper-style encoders
\citep{radford2023whisper}, so the 32-layer encoder always processes 1500 frames,
although course queries last 2.2--10.4\,s (mean 5.2\,s on practice, 5.5\,s on dev).

\textbf{Question.} \emph{How much of ASR latency in this pipeline is redundant
computation, how much of it can be removed exactly (KV cache) or approximately
(encoder trimming), and what does the approximate removal cost in transcript and
retrieval quality?}

\textbf{Contributions.}
\begin{itemize}
  \item A KV-cache decoding path (V2) that reuses the supplied operations and our
        kernels, including a RoPE kernel with a position offset, whose tokens and
        probabilities match the recompute baseline bit for bit.
  \item Encoder trimming (V3) with a margin sweep and a 10\,s bucket that never cuts
        real audio, giving a latency--quality curve rather than a single point.
  \item A cost model written before any measurement (\S\ref{sec:costmodel}), compared
        against the measured speed-ups, and a mechanism-level account of where the
        remaining time goes once the redundancy is gone.
  \item An evaluation with one pinned GPU type, separated timing boundaries, and
        quality measured against retrieval relevance labels, not only transcript agreement.
\end{itemize}
Answers to the four required HW1 guide questions are in Appendix~\ref{app:hw1}.

% ============================================================
\section{Background and baseline}
\label{sec:background}

\textbf{Model shapes.} The audio encoder has 32 layers with hidden size 1280,
20 heads of dimension 64, an FFN of 5120 and non-causal attention over $T=1500$
positions (3000 mel frames after a stride-2 convolution). A projector pools 4
encoder positions into one audio token. The text decoder has 28 layers, hidden size
2048, 16 query heads and 4 KV heads of dimension 128 (grouped-query attention),
an FFN of 6144 and a 59{,}264-token vocabulary; it holds $\approx$1.35\,B parameters,
or 5.4\,GB in float32. The prompt is a few text tokens plus one audio token per 4
encoder positions of \emph{real} audio, so it is 42--144 positions on practice
(median 77); on the supplied 3.5\,s test clip it is 59, and the 13 generated tokens
grow the sequence to 71.

\textbf{HW1 kernels.} We implemented the four task kernels the model runs on:
RoPE (Task~1, \citealp{su2021roformer}), warp/block reductions with RMSNorm,
LayerNorm and a single-pass online softmax (Task~2, \citealp{milakov2018online}),
a shared-memory tiled GEMM with a fused GELU epilogue (Task~3), and a tiled
attention kernel with an online softmax over key tiles (Task~4,
\citealp{vaswani2017attention,dao2022flashattention}). All are hand-written CUDA; no
cuBLAS, cuDNN or CUTLASS is used on the CUDA path.

\textbf{Baselines.} V0 is the supplied model on PyTorch (\texttt{ASR\_CUDA=none}).
V1 is the same model with every operation on our kernels (\texttt{ASR\_CUDA=all});
it is the baseline for the investigation. On an NVIDIA GeForce RTX 2080 Ti (sm\_75,
68 SMs, 616\,GB/s), \texttt{make test} passes all four tasks with maximum errors of
$1.9{\times}10^{-7}$ (RoPE), $5.4{\times}10^{-7}$ (reductions), $3.6{\times}10^{-5}$
(GEMM, $K=6144$) and $5.2{\times}10^{-7}$ (attention). Our GEMM reaches 1.38\,TFLOP/s
on encoder fc1 against 13.4\,TFLOP/s for cuBLAS. \texttt{./run.sh --profile} reports
3290\,ms per transcription on our kernels against 877\,ms on PyTorch; on our kernels
linear layers take 86.7\% of kernel time and attention 12.8\%, and every other
operation together under 1\%. The GEMM gap therefore sets the scale of every
latency in this report (Appendix~\ref{app:hw1}, Task~3). Full logs are in
Appendix~\ref{app:logs}.

% ============================================================
\section{Design}
\label{sec:design}

\begin{table}[t]
\centering
\small
\caption{Variants. Every variant runs the supplied \texttt{ops.*} functions, so
\texttt{ASR\_CUDA} selects kernels as in V0/V1; \texttt{model.py} is unchanged.}
\label{tab:variants}
\setlength{\tabcolsep}{4pt}
\begin{tabular}{lll}
\toprule
ID & Configuration & Role \\
\midrule
V0 & supplied model, PyTorch path & supplied reference \\
V1 & supplied model, our kernels & investigation baseline \\
V2 & V1 + KV-cache decoding & exact optimisation \\
V3-$m$ & V2 + encoder input cut to real length + margin $m$ & approximate \\
V3-b10 & V2 + encoder input fixed at 10\,s (never cuts audio) & approximate \\
\bottomrule
\end{tabular}
\end{table}

\textbf{V2: KV cache.} Each decoder layer gets preallocated key and value caches of
shape $[1, S_{\max}, 4, 128]$ with $S_{\max} = P + N_{\max}$ for prompt length $P$
and token budget $N_{\max}$ ($\approx$28\,MB for the longest query;
cf.\ \citealp{kwon2023pagedattention,pope2023efficiently}). Prefill runs the decoder
once over the prompt with causal attention and writes rows $[0, P)$. Each later step
runs one position: $q/k/v$ projections with $M=1$, RoPE at position $P+t-1$, a write
into cache row $P+t-1$, and attention of one query against the cache prefix. The
attention kernel's causal shift ($\mathit{seq}_k - \mathit{seq}_q$) already gives the
correct window, and a prefix slice of a batch-1 cache is contiguous, so no copy is
needed. Greedy argmax and the EOS set are those of the supplied loop. Because every
output row of our tiled GEMM accumulates over $K$ in the same order regardless of
$M$, and attention is per query, we \emph{hypothesise} bit-identical probabilities;
\S\ref{sec:results} tests it.

\textbf{RoPE offset.} The supplied RoPE kernel rotates positions $0..S-1$. We
measured two ways to rotate one token at position $p$: (a) a new entry point
\texttt{launch\_rope\_at} in \texttt{task1\_rope.cu} that shares a device helper with
the original kernel and takes \texttt{pos\_offset} (existing names and signatures
untouched); (b) a fallback that runs the unchanged \texttt{launch\_rope} over a
padded $[1, p{+}1, \dots]$ scratch tensor and keeps the last row. Variant (b) needs no
kernel change but its cost grows with $p$. Both give bit-identical outputs (unit test).
On 20 practice queries (V2) they gave identical transcripts and 38.3 vs 38.6\,ms per
token: at $p<160$ the padded work is negligible next to the decoder GEMMs. Every
reported number uses (a).

\textbf{V3: encoder trimming.} From the feature mask we take the real mel length
$L$ and keep $T'_{\mathrm{mel}} = \min(3000, \mathrm{roundup}_8(L + 100m))$ frames for
margin $m \in \{0, 1, 2, 5\}$\,s (rounding to 8 matches conv stride 2 $\times$ pool 4).
The bucket setting b10 uses $T'_{\mathrm{mel}} = \max(1000, \mathrm{roundup}_8(L))$,
a fixed shape for every query up to 10\,s that still never cuts audio (the longest
query is 10.4\,s). In every setting $T' \ge L$, so the number of audio tokens and the
prompt length are unchanged; a unit test asserts this. The approximation is that real
frames no longer attend to the removed padding frames in the non-causal encoder,
which is the mechanism that can change transcripts.

\textbf{Timing structure.} We time CPU feature extraction (\texttt{feat\_ms}),
encoder, prefill (which produces the first token) and the remaining decode steps
separately with CUDA events. \texttt{ms\_per\_token} is decode time after the first
token divided by the number of such steps, so prefill cost and per-step cost are not
mixed. V1 is timed by a driver that calls the supplied \texttt{encode\_audio} and
\texttt{next\_token\_probs} in the supplied control flow, with events in between; a
test checks it reproduces \texttt{generate()} exactly.

% ============================================================
\section{Cost model (written before measurement)}
\label{sec:costmodel}

\textbf{Decode floor.} After the KV cache, each step multiplies a single activation
row by every decoder weight matrix, so it must stream all weights from DRAM once:
$1.35\,\text{B} \times 4\,\text{B} \approx 5.4$\,GB, or $\approx$8.8\,ms per token at
the RTX 2080 Ti's 616\,GB/s; the lm\_head adds 0.49\,GB (0.8\,ms). This is a roofline
bound \citep{williams2009roofline}: at $M=1$ the arithmetic intensity is
$\approx$0.5\,FLOP/byte, far below the ridge point.

\textbf{V1 recompute.} V1 runs the full decoder over prefixes of 59 to 70 positions,
$\sum_{S=59}^{70} S = 774$ position-passes for its 12 steps after the first token,
against $59 + 12 = 71$ for V2. Each position costs $2 \times 1.35$\,GFLOP,
so a V1 step is 159--189\,GFLOP. Our tiled GEMM also re-reads each weight matrix once
per row tile, $\lceil S/\mathrm{TILE} \rceil = 4$ times for TILE${}=16$, so a V1 step
moves $\gtrsim$21.6\,GB. We predict V1 step time
$\approx \max(\text{FLOP}/F_{\mathrm{ach}},\ 4 \times 8.8\,\text{ms})$ with $F_{\mathrm{ach}}$
the GFLOP/s our GEMM reaches on decoder shapes in \texttt{make test}.

\textbf{Encoder trimming.} Per layer, linear layers cost $\approx 39.3\,\text{M}\cdot T$
FLOP and attention $5120\,T^2$ FLOP, so at $T=1500$ attention is 16\% of encoder
FLOPs (2.26\,TFLOP in total). Table~\ref{tab:costmodel} gives the predicted encoder
ratio $E(T')/E(1500)$ for the median 5.0\,s query.

\begin{table}[t]
\centering
\small
\caption{Predictions fixed before measurement (median 5.0\,s query, $L=500$ mel
frames; RTX 2080 Ti, 616\,GB/s). Encoder ratios count FLOPs of the 32 transformer
layers only (convolutions and projector excluded). \emph{Measured}: median ratio of
encoder ms and median ms per token on practice ($n=350$); m1 and m5 were not run.}
\label{tab:costmodel}
\setlength{\tabcolsep}{4.5pt}
\begin{tabular}{lrrrrrr}
\toprule
Setting & $T'_{\mathrm{mel}}$ & $T'$ & linear & attention & $E(T')/E(1500)$ & measured \\
\midrule
full (V2) & 3000 & 1500 & 1.000 & 1.000 & 1.000 & 1.000 \\
m0  & 504  & 252 & 0.168 & 0.028 & 0.145 & 0.136 \\
m1  & 600  & 300 & 0.200 & 0.040 & 0.174 & -- \\
m2  & 704  & 352 & 0.235 & 0.055 & 0.205 & 0.192 \\
m5  & 1000 & 500 & 0.333 & 0.111 & 0.297 & -- \\
b10 & 1000 & 500 & 0.333 & 0.111 & 0.297 & 0.288 \\
\midrule
\multicolumn{5}{l}{V2 decode, ms/token floor (5.4\,GB / 616\,GB/s; +0.8\,ms lm\_head)} & 8.8 & 38.4 \\
\multicolumn{5}{l}{V1 step, weight-traffic bound ($4 \times 5.4$\,GB at TILE${}=16$)} & 35 & 165.6 \\
\bottomrule
\end{tabular}
\end{table}

Combining terms, the predicted V2 latency is $E + D_P + \sum_t(\mathrm{step}_t + H)$
and V3 replaces $E$ by $E \cdot E(T')/E(1500)$. The predictions in
Table~\ref{tab:costmodel} were committed (\texttt{dc44796}) before any variant was timed.

% ============================================================
\section{Experimental setup}
\label{sec:setup}

\textbf{Hardware and boundaries.} Every timing comes from one pinned GPU type
(\texttt{--gres=gpu:nvidia\_geforce\_rtx\_2080\_ti:1}; \texttt{nvidia-smi -L}:
``NVIDIA GeForce RTX 2080 Ti''; CUDA 12.8; all variant runs on one node). The ASR
timing boundary is a warm model, from a 16\,kHz float array in to a token list out;
model load (26--113\,s, dominated by the shared file system) is excluded. Retrieval
scoring runs as a CPU-only job (Intel Xeon E5-2620 v4, 4 cores) from cached
transcripts, which the course allows; it does not contribute to any reported
latency, and every variant is scored on the same machine and settings.

\textbf{Data and protocol.} Development used the practice split (350 queries); the
final tables use the dev split (125 queries), run once after the design was frozen at
\texttt{f9f1ab3}. Each run discards 5 warm-up queries, times every query once, then
repeats a fixed 50-query subset to estimate run-to-run variance (median per-query
difference 0.1--0.8\%).

\textbf{Metrics.} Cost: ASR latency mean and p95, \texttt{prefill\_ms},
\texttt{ms\_per\_token}, encoder and decoder ms, and peak GPU memory. Quality:
transcript exact-match rate and word error rate against V1 (there are no reference
transcripts), bit-identity of token probabilities for V2 (unit test), and Recall@10,
MRR@10 and NDCG@10 against the relevance labels with the supplied enc\_a encoder and
exact search, with a paired bootstrap interval on the per-query difference from V1.
We planned an LLM judge of generated answers on the queries whose transcript changed;
ELM access did not arrive in time, so answer quality is not measured
(\S\ref{sec:conclusion}).

% ============================================================
\section{Results}
\label{sec:results}

\begin{table}[t]
\centering
\caption{ASR cost and quality per variant, dev split (125 queries), RTX 2080 Ti.
R@10 brackets give the paired difference from V1 with a 95\% bootstrap interval.
Generated by \texttt{summarise.py}.}
\label{tab:main}
\IfFileExists{tables/asr_dev.tex}{\resizebox{\linewidth}{!}{\begin{tabular}{lrrrrrrrr}
\toprule
Variant & ASR ms (mean) & p95 & enc ms & dec ms & ms/token & match \% & WER & R@10 \\
\midrule
V1 & 4971.3 & 7631.7 & 1761.1 & 3186.2 & 186.7 & 100.0 & 0.000 & 0.475 \\
V2 & 2494.5 & 2845.8 & 1709.0 & 754.6 & 38.0 & 100.0 & 0.000 & 0.475 (+0.000 [+0.000, +0.000]) \\
V3 b10 & 1255.8 & 1597.8 & 484.7 & 753.8 & 38.2 & 96.0 & 0.003 & 0.475 (+0.000 [+0.000, +0.000]) \\
V3 m0 & 1036.4 & 1476.4 & 260.4 & 748.4 & 37.8 & 91.2 & 0.008 & 0.475 (+0.000 [-0.012, +0.012]) \\
V3 m2 & 1128.6 & 1580.0 & 358.7 & 748.1 & 37.8 & 96.0 & 0.003 & 0.475 (+0.000 [+0.000, +0.000]) \\
\bottomrule
\end{tabular}
% split=dev, n per variant=[125, 125, 125, 125, 125], GPU=NVIDIA GeForce RTX 2080 Ti
}}{\TBD{Table 1: \texttt{tables/asr\_dev.tex} not generated yet}}
\end{table}

\begin{figure}[t]
\centering
\figorTBD{figures/breakdown_dev.pdf}{0.48\linewidth}{Fig.\ 1: latency breakdown}\hfill
\figorTBD{figures/tradeoff_dev.pdf}{0.48\linewidth}{Fig.\ 2: latency vs Recall@10}
\caption{Left: mean ASR latency split into CPU features, encoder and decoder.
Right: latency against Recall@10 over trim settings. Dev, $n=125$, RTX 2080 Ti.}
\label{fig:breakdown}
\end{figure}

\textbf{KV cache (exact).} V2 cuts mean ASR latency from 4971 to 2495\,ms (2.0$\times$)
and decoder time from 3186 to 755\,ms (4.2$\times$); the median step falls from 186.8
to 38.1\,ms. All 125 transcripts are identical to V1, and the unit test confirms the
token probabilities are bit-identical, as hypothesised in \S\ref{sec:design}. The
cache costs at most 26\,MB against 8.6\,GiB of weights, so peak memory is unchanged.
A V2 step streams the 5.9\,GB of decoder and lm\_head weights in 38.1\,ms, i.e.\
155\,GB/s or 25\% of peak bandwidth: 4.3$\times$ above the 8.8\,ms floor. Prefill
(the whole prompt plus the first token) takes 167--179\,ms in every variant. With
decoding cheap, the unchanged 1.71\,s encoder becomes 69\% of V2's latency, which is
what V3 attacks.

\textbf{Encoder trimming (approximate).} Trimming to the real length (m0) cuts the
median encoder time from 1717 to 266\,ms; with a 2\,s margin it is 359\,ms and with
the 10\,s bucket 486\,ms. Mean ASR latency becomes 1036, 1129 and 1256\,ms, i.e.\
4.8$\times$, 4.4$\times$ and 4.0$\times$ faster than V1. The cost is small: 11, 5 and
5 of 125 transcripts change (WER against V1 0.8\%, 0.3\% and 0.3\%). Recall@10 is
unchanged for every query under m2 and b10, and under m0 the change is $+0.000$
$[-0.012, +0.012]$; MRR@10 moves by at most $+0.004$. No setting changes retrieval
measurably, and the practice split agrees (4.7$\times$ for m0; every $\Delta$Recall@10
interval within $\pm0.01$; Appendix~\ref{app:practice}). That b10 still changes
transcripts shows the mechanism: no real audio is cut in any setting, so every change
comes from real frames no longer attending to padding, and how much padding remains
matters less than whether it is the full 30\,s. After trimming, the decoder (prefill
plus steps, $\approx$910\,ms) is 88\% of V3-m0's latency.

\begin{figure}[t]
\centering
\figorTBD{figures/speedup_dev.pdf}{0.6\linewidth}{Fig.\ 3: speed-up vs clip duration}
\caption{Per-query ASR speed-up over V1 against audio duration. Dev, $n=125$,
RTX 2080 Ti.}
\label{fig:speedup}
\end{figure}

\textbf{Prediction vs measurement.} On practice, whose median query the model was
written for, the encoder ratios are 3--7\% below the FLOP model in every setting; on
dev, b10 (the one setting whose shape does not depend on clip length) gives 0.283,
5\% below. The model weighs attention FLOPs like GEMM FLOPs, but our
attention kernel runs at 959\,GFLOP/s against 1380--1860 for our GEMM
(\texttt{make test}), so the quadratic term is a larger share of \emph{time} at
$T=1500$ than of FLOPs, and removing it saves more than predicted. The practice V1 step
(165.6\,ms) is far above the 35\,ms weight-traffic bound: V1 is compute-bound. Its
FLOP term, at the 1.27\,TFLOP/s our GEMM reaches on decoder shapes, predicts 180\,ms;
the 9\% overestimate fits our GEMM running faster with the 77--90 rows of a real
prefix than with the 64 rows of the test shapes. The V2 decode floor is the one
prediction far from the measurement (8.8 vs 38.4\,ms); \S\ref{sec:discussion}
explains the gap.

% ============================================================
\section{Discussion}
\label{sec:discussion}

\textbf{Decode becomes GEMV.} Once prefix recomputation is gone, every decoder
matrix multiply has $M=1$. Our tiled GEMM launches TILE$\times$TILE thread blocks, of
which only one row does useful work at $M=1$, so (TILE${}-1$)/TILE of the threads
idle; every block still loads its weight tile and synchronises twice per $K$ tile,
so the kernel is limited by memory latency, not FLOPs. \texttt{make test} shows the
effect in isolation: lm\_head at $M=1$ runs at 123\,GFLOP/s, i.e.\ 246\,GB/s or
40\% of peak bandwidth, against cuBLAS's 291\,GFLOP/s. Across a whole step, with
smaller matrices and per-layer launches, we reach 25\%. A GEMV kernel (one warp per
output column, coalesced float4 loads, no shared-memory tile) is the direct fix;
this is HW1 Task~3's TILE trade-off seen from the decode side.

\textbf{Attention at $\mathit{seq}_q=1$.} The attention kernel launches one block per
(batch, head, query tile); with one query that is 16 blocks on a 68-SM GPU, and only
one of ATTN\_WARPS warps has a query row. With 77--160 keys this is cheap in absolute
terms, but it leaves most SMs idle during the attention launch of every layer.

\textbf{Considered but not built.} \emph{CUDA Graphs} for the decode step would remove
launch overhead from $\approx$300 kernel launches per token, but the sequence length
changes every step, so graphs would need re-capture or padded shapes, and the risk
was high for the time available. \emph{Split-K attention} for $\mathit{seq}_q=1$ would
fix the 16-block occupancy problem, but attention is only 12.8\% of kernel time even
in V1, and less once steps are weight-bandwidth-bound. \emph{FP16 weights} would halve the
bandwidth floor to $\approx$4.4\,ms/token, but are inexact, break the bit-identical
claim, and change the float32 kernel inputs the HW1 rules fix. \emph{VAD-based
trimming} per clip would need an extra model in the pipeline, while the real-length
floor already never cuts audio.

% ============================================================
\section{Limitations and conclusion}
\label{sec:conclusion}

\textbf{Limitations.} All timings come from one GPU type and batch size 1; the
bandwidth argument predicts different ratios on GPUs with a different
FLOP-to-bandwidth balance. Without reference transcripts, WER is measured against V1,
so it measures change rather than error. The dev split has 125 queries, which limits
how small a retrieval difference can be resolved: the widest 95\% interval on
$\Delta$Recall@10 is $\pm0.012$ (dev, m0), and $\pm0.01$ on practice ($n=350$). The ELM judge was not run, so
we show that trimming does not change \emph{retrieval}, not that it never changes an
answer; the 4--9\% of changed transcripts are where answers could differ.

\textbf{Conclusion.} About 79\% of the baseline ASR time was redundant work: the KV
cache removes the prefix recomputation exactly (2.0$\times$, bit-identical), and
trimming the encoder to the real audio removes the padding approximately (4.8$\times$
in total) at no measurable retrieval cost. What remains is a decoder that streams
5.9\,GB of weights per token at a quarter of peak bandwidth, so a GEMV kernel for
$M=1$ is the next step, not more redundancy removal.

% References are excluded from the main-body page limit.
\bibliography{references}
\bibliographystyle{iclr2027_conference}

% ============================================================
\appendix

\section{HW1 kernel answers}
\label{app:hw1}

One question per guide, quoted verbatim from \texttt{hw1-asr/docs/}. All numbers come
from the pinned RTX 2080 Ti run in Appendix~\ref{app:logs}.

\subsection{Task 1 (RoPE), Question 1}
\begin{quote}
``For the audio encoder case, how many bytes does the kernel read and write in total?
Which fraction of your GPU's peak bandwidth did you reach?''
\end{quote}

The encoder uses partial rotation: only the first 32 of 64 dimensions per head
are rotated, and the kernel touches only those. $q$ and $k$ each have
$1500 \times 20 = 30{,}000$ rows, so it reads $60{,}000 \times 32 \times 4$\,B
$= 7.68$\,MB and writes the same amount: 15.36\,MB in total, plus 64\,B of
\texttt{inv\_freq} (cached). The other 32 dimensions are never loaded.

\texttt{make test1} measures 0.036\,ms, i.e.\ 424\,GB/s, which is 69\% of the
RTX 2080 Ti's 616\,GB/s. The small decoder case (0.006\,ms) reaches only 38\%:
at that size, launch and tail latency dominate.

\subsection{Task 2 (reductions), Question 3}
\begin{quote}
``How many times does your softmax read each element of \texttt{x} from global memory?
How many times does the textbook version (max, then sum, then write) read it?''
\end{quote}

Twice. Pass 1 reads each element once and keeps a per-thread running
(max, sum) pair, rescaling the sum when the max grows; the pairs are then
combined by a block reduction. Pass 2 reads each element again to write
$e^{x-\max}/\mathrm{sum}$. Traffic: 2 reads + 1 write per element.

The textbook version reads three times (max, sum, normalise-and-write): 3 reads
+ 1 write, so the online version moves 25\% less data. The test counts 1 read + 1
write, so its 194\,GB/s on the $64 \times 4096$ case is really
$\approx$290\,GB/s (47\% of peak). The single 59{,}264-wide vocabulary row runs
at 3.7\,GB/s: it is one block on one of 68 SMs, so it is latency-bound, not
bandwidth-bound (0.13\,ms per token).

\subsection{Task 3 (GEMM), Question 2}
\begin{quote}
``Your kernel reaches only a fraction of cuBLAS's speed on the large cases. Name one
thing cuBLAS does differently that explains part of the gap. You do not have to
implement it.''
\end{quote}

On encoder fc1 ($1500 \times 5120 \times 1280$, + GELU) our kernel reaches
1383\,GFLOP/s and cuBLAS 13{,}371\,GFLOP/s, about the GPU's FP32 peak: we are at
10\%. Across the large cases we reach 10--28\% of cuBLAS.

\emph{Register blocking.} Each of our threads computes one output, so every
fused multiply-add needs two shared-memory loads (\texttt{xs[ty][kk]},
\texttt{ws[tx][kk]}), and shared-memory load throughput per SM is well below
its FMA throughput, so the kernel is bound by shared-memory loads and
\texttt{\_\_syncthreads}, not arithmetic. cuBLAS gives each thread a block of outputs, e.g.\
$8 \times 8$, held in registers. It loads 8 values of $x$ and 8 of $w$ per $k$
and does 64 FMAs with them, 4 FMAs per load instead of 0.5, so FMA throughput
becomes the limit. It also uses vectorised loads and double-buffers the next
tile while computing on the current one.

\subsection{Task 4 (attention), Question 1}
\begin{quote}
``How much memory would the full score matrix of one encoder layer take (20 heads,
1500 x 1500, float32)? How much shared memory does one block of your kernel use
instead?''
\end{quote}

$20 \times 1500 \times 1500 \times 4$\,B $= 180$\,MB (172\,MiB) per layer;
across 32 layers it would be written and read 32 times per transcription.

One block holds a key tile of 32 keys (padded to $\mathit{head\_dim}+1$ to
avoid bank conflicts), a value tile of 32 keys and ATTN\_WARPS${}=4$ query rows:
$4 \times (32 \cdot 65 + 32 \cdot 64 + 4 \cdot 64) = 17{,}536$\,B (17.1\,KiB) at
$\mathit{head\_dim}=64$, independent of sequence length. This is the dynamic
shared-memory size the launcher computes. Scores never leave registers: each
lane holds one score per tile, and the online softmax folds it into the running
output.

\section{Full test and profile logs}
\label{app:logs}

Environment: NVIDIA GeForce RTX 2080 Ti (sm\_75), CUDA 12.8, Informatics teaching
cluster (\texttt{Teaching} partition). All three logs are from one session on one node.

\subsection{\texttt{make test}}
Each line is one test case: the maximum absolute error against the CPU reference, the
kernel time, and its bandwidth (Tasks 1--2) or GFLOP/s (Tasks 3--4). For GEMM the
naive kernel and cuBLAS are shown for comparison; cuBLAS is used only by the test.

\logorTBD{logs/make_test.txt}{full \texttt{make test} output}

\subsection{\texttt{./run.sh --profile} (V1, our kernels)}
Three timed transcriptions of the supplied 3.5\,s test clip after warm-up, then a
per-operation breakdown of one transcription measured with CUDA events. Model load is
not included in the per-transcription time; it depends on the shared file system.

\logorTBD{logs/profile_cuda.txt}{\texttt{./run.sh --profile} output}

\subsection{\texttt{./run.sh --cuda none --profile} (V0, PyTorch)}
The same run with every operation on the PyTorch reference path (cuBLAS for linear
layers), on the same node. The gap to V1 is almost entirely in \texttt{linear}.

\logorTBD{logs/profile_torch.txt}{\texttt{./run.sh --cuda none --profile} output}

\subsection{Practice split results}
\label{app:practice}
\IfFileExists{tables/asr_practice.tex}{{\centering\resizebox{\linewidth}{!}{\begin{tabular}{lrrrrrrrl}
\toprule
Variant & ASR & p95 & enc & dec & step & chg. & WER & R@10; $\Delta$ [95\% CI] \\
\midrule
V1 & 4580 & 7867 & 1704 & 2858 & 169.9 & 0 & 0.0 & 0.480 \\
V2 & 2453 & 2921 & 1698 & 734 & 38.5 & 0 & 0.0 & identical \\
V3-m0 & 991 & 1570 & 245 & 729 & 38.2 & 22 & 0.8 & $+0.004$ [$-0.006$, $+0.014$] \\
V3-m2 & 1092 & 1684 & 342 & 732 & 38.4 & 19 & 0.4 & $-0.002$ [$-0.009$, $+0.002$] \\
V3-b10 & 1238 & 1689 & 488 & 732 & 38.4 & 18 & 0.5 & $+0.004$ [$+0.000$, $+0.010$] \\
\bottomrule
\end{tabular}
% split=practice, n per variant=[350, 350, 350, 350, 350], GPU=NVIDIA GeForce RTX 2080 Ti
}\par}}{\TBD{\texttt{tables/asr\_practice.tex} not generated yet}}

\section{Reproducibility}
\label{app:repro}

Variant runs: commit \texttt{f9f1ab3}. All commands run from the repository root;
\texttt{scripts/*.sbatch} wrap them for Slurm on the pinned GPU type.
\begin{small}
\begin{verbatim}
make -C hw1-asr && make -C hw1-asr test
(cd hw1-asr && ./run.sh --profile)
# one job per variant (v1, v2; v3 with --margin 0|2 or --bucket 10)
sbatch scripts/bench.sbatch results/asr dev_v3_m0 --split dev \
    --variant v3 --margin 0 --repeat-subset 50
sbatch scripts/retrieval.sbatch dev v1 v2 v3_m0 v3_m2 v3_b10   # CPU only
python rag/extensions/analysis/summarise.py --split dev
\end{verbatim}
\end{small}
\texttt{scripts/check\_kernels.sbatch} runs both profiles and the extension test
suite (48 tests, including V2 bit-identity and the two RoPE paths) on the same GPU type.

The supplementary zip (\texttt{scripts/make\_submission\_zip.sh}) contains the five
HW1 task files, \texttt{rag/extensions/} (KV-cache and trimming variants, timing
harness, variant ASR worker, analysis), a README separating our code from supplied
code, and the measured CSV and text results. It contains no data, model weights,
caches or credentials.

% The AI declaration / student information page is added separately as the
% final page (template: \newpage, \section*{AI declaration}, group information).
% It is intentionally not part of this anonymous skeleton.

\end{document}
